\section{Tool Integration via MCP}

\label{sec:tools}

\subsection{Model Context Protocol}

The Model Context Protocol (MCP)~\cite{anthropic2024mcp} provides a standardized interface for LLMs to discover and invoke external tools. Hanzo Chat implements an MCP client that connects to multiple MCP servers simultaneously.

\subsection{Tool Categories}

Hanzo Chat ships with 260+ tools organized into categories:

\begin{table}[H]
\centering
\small
\begin{tabular}{lcc}
\toprule
\textbf{Category} & \textbf{Tools} & \textbf{Example} \\
\midrule
Code execution & 15 & Python, Node.js, Bash \\
File operations & 22 & Read, write, search, glob \\
Web interaction & 18 & Fetch, search, browse \\
Database & 12 & SQL, vector, key-value \\
API integration & 45 & REST, GraphQL, gRPC \\
DevOps & 28 & Docker, K8s, CI/CD \\
Data analysis & 20 & Pandas, plotting, stats \\
Communication & 15 & Email, Slack, Discord \\
Version control & 18 & Git, GitHub, PRs \\
Cloud services & 35 & AWS, GCP, DO \\
AI/ML & 22 & Training, inference, eval \\
System & 30 & OS, process, network \\
\bottomrule
\end{tabular}
\caption{MCP tool categories and counts.}
\label{tab:tools}
\end{table}

\subsection{Tool Discovery and Selection}

Not all 260+ tools are included in every prompt. Hanzo Chat uses a two-stage tool selection process:

\begin{algorithm}[H]
\caption{Dynamic Tool Selection}
\label{alg:tools}
\begin{algorithmic}[1]
\Require Query $m_t$, query embedding $\bm{e}_t$, all tools $\mathcal{T}_{\text{all}}$
\State \Comment{Stage 1: Embedding-based retrieval}
\State $\mathcal{T}_{\text{cand}} \gets \text{Top-}20(\cos(\bm{e}_t, \bm{e}_{\text{tool}}))$
\State \Comment{Stage 2: LLM-based filtering}
\State $\mathcal{T}_{\text{sel}} \gets \text{LLM.Filter}(m_t, \mathcal{T}_{\text{cand}})$
\State \Comment{Always include core tools}
\State $\mathcal{T}_{\text{final}} \gets \mathcal{T}_{\text{sel}} \cup \mathcal{T}_{\text{core}}$
\State \Return $\mathcal{T}_{\text{final}}$ \Comment{Typically 8--15 tools}
\end{algorithmic}
\end{algorithm}

This reduces the tool definition overhead from $\sim$50K tokens (all tools) to $\sim$4K tokens (selected subset) while maintaining 97\% recall on tool-use benchmarks.

\subsection{Tool Execution Sandboxing}

All tool executions run in isolated environments:

\begin{itemize}[leftmargin=1.1em]
    \item \textbf{Code execution}: Docker containers with resource limits (CPU: 2 cores, RAM: 4GB, time: 60s, no network unless explicitly granted).
    \item \textbf{File operations}: Scoped to project directories with read/write permissions managed per user.
    \item \textbf{Web requests}: Rate-limited (10 req/min), filtered for malicious URLs, response size capped at 1MB.
    \item \textbf{API calls}: OAuth credentials stored in encrypted vault, never exposed to the LLM.
\end{itemize}

\subsection{Multi-Step Tool Chains}

Hanzo Chat supports iterative tool use where the model can invoke multiple tools in sequence, using the output of one as input to the next. The execution loop (Algorithm~\ref{alg:pipeline}, lines 7--10) continues until the model produces a final text response without tool calls.

We observe that complex tasks often require 3--7 tool calls. To prevent infinite loops, we impose a configurable maximum (default: 25 tool calls per turn).

\begin{table}[H]
\centering
\small
\begin{tabular}{lcc}
\toprule
\textbf{Task Type} & \textbf{Avg. Tool Calls} & \textbf{Success Rate} \\
\midrule
Simple file read & 1.2 & 99.1\% \\
Code debugging & 3.8 & 87.4\% \\
Data analysis & 5.1 & 82.3\% \\
Multi-file refactor & 7.3 & 78.9\% \\
Full-stack feature & 12.4 & 71.2\% \\
\bottomrule
\end{tabular}
\caption{Tool call statistics by task complexity.}
\label{tab:tool_stats}
\end{table}

